\section{The Son of Man in His Majesty}\label{sec:majesty}

As the liturgical year draws to its close the Church lifts her eyes to its end. On the Last Sunday after Pentecost the 1962 Missal reads the discourse that Christ gave his disciples ``sitting on mount Olivet,'' when they asked him what should be the sign of his coming and of the consummation of the world (Matt 24:3). He speaks of the abomination of desolation standing in the holy place, of a tribulation such as has not been from the beginning of the world, of false christs and false prophets, of the sun darkened and the stars falling from heaven; and then of himself:

\begin{quote}
And then shall appear the sign of the Son of man in heaven. And then shall all tribes of the earth mourn: and they shall see the Son of man coming in the clouds of heaven with much power and majesty. And he shall send his angels with a trumpet and a great voice: and they shall gather together his elect from the four winds, from the farthest parts of the heavens to the utmost bounds of them. (Matt 24:30--31)
\end{quote}

The Gospel of that Sunday ends on a sentence that gathers the whole year into its hand: ``Heaven and earth shall pass: but my words shall not pass'' (Matt 24:35). A week later the year begins again, and the first Sunday of Advent reads in the same Missal the same prophecy as Luke records it, with the sun and moon and stars, the distress of nations, the roaring of the sea, and men ``withering away for fear and expectation of what shall come upon the whole world.'' Then, into the midst of the terror, a command: ``But when these things begin to come to pass, look up and lift up your heads, because your redemption is at hand'' (Luke 21:25--28). The Church ends one year and begins the next with the same coming before her eyes, read once in fear and once in hope.

In the Postconciliar Missal the year closes on the Solemnity of Christ the King, and in the year of Matthew (Year~A) its Gospel is the scene with which the discourse on Olivet ends, the last thing Christ taught before he said to his disciples that after two days the Son of man would be delivered up to be crucified (Matt 26:1--2):

\begin{quote}
And when the Son of man shall come in his majesty, and all the angels with him, then shall he sit upon the seat of his majesty. And all nations shall be gathered together before him: and he shall separate them one from another, as the shepherd separateth the sheep from the goats: and he shall set the sheep on his right hand, but the goats on his left. Then shall the king say to them that shall be on his right hand: Come, ye blessed of my Father, possess you the kingdom prepared for you from the foundation of the world. (Matt 25:31--34)
\end{quote}

Then follow the six works: the hungry fed, the thirsty given drink, the stranger taken in, the naked clothed, the sick visited, the prisoner sought out; and the word that turns every one of them toward Christ himself: ``Amen I say to you, as long as you did it to one of these my least brethren, you did it to me'' (Matt 25:40). To those on the left he says: ``Depart from me, you cursed, into everlasting fire, which was prepared for the devil and his angels'' (25:41). And the scene closes on a single sentence with two halves of equal weight: ``And these shall go into everlasting punishment: but the just, into life everlasting'' (25:46).

The Church has never treated this as one teaching among others. She confesses it every Sunday in the Creed, which says of the Son that ``he will come again in glory to judge the living and the dead'';\endnote{The Creed in the English quoted at \emph{Spe salvi} 41.} and Benedict~XVI observes that from the earliest times this truth of the Creed ``has influenced Christians in their daily living as a criterion by which to order their present life, as a summons to their conscience, and at the same time as hope in God's justice.''\endnote{Benedict~XVI, encyclical letter \emph{Spe salvi} (30 November 2007), 41, official English.} Order, summons, and hope: those three words will serve as well as any to describe what the doctrine of judgment is for.

\subsection{The Church's own meditation}

The Church has also given her children the words in which to meditate on that day. On the second of November, when she commemorates all the faithful departed and every priest may offer the Holy Sacrifice three times for the dead, the Mass of Requiem in the 1962 Missal sets between the Epistle and the Gospel the sequence in which the Church herself meditates on the last judgment. It does not begin with an argument. It begins with a day:

\begin{quote}
\latin{Dies irae, dies illa,\\
\quad Solvet saeclum in favilla:\\
\quad Teste David cum Sibylla.\\
Quantus tremor est futurus,\\
\quad Quando judex est venturus,\\
\quad Cuncta stricte discussurus!}\endnote{\emph{Missale Romanum} (Tournai: Desclée, Lefebvre et Soc., 1889), \emph{Missae pro Defunctis}, Sequence, p.~92*. Stress accents omitted.}

\medskip
Day of Wrath! O day of mourning!\\
See fulfill'd the prophets' warning!\\
Heav'n and earth in ashes burning!\\
Oh, what fear man's bosom rendeth\\
When from Heav'n the Judge descendeth,\\
On Whose sentence all dependeth!\endnote{W.~J. Irons, ``Day of Wrath! O day of mourning,'' stanza 1, in \emph{Hymns Ancient and Modern}, complete edition (London: printed for the Proprietors by William Clowes and Sons, n.d., catalogued 1890), hymn 398, p.~307.}
\end{quote}

The sequence moves as the soul moves when it truly considers the judgment. It begins with the cosmos---the world dissolving into ash, the trumpet sounding through the sepulchres of every land, death and nature struck dumb as the creature rises to answer its Judge. It narrows to a book in which everything is written, and to a Judge before whom whatever is hidden shall appear. Then, without warning, it falls from the third person into the first: \latin{Quid sum miser tunc dicturus?} What shall I, wretched one, say then? From that line to the end it is no longer a description of the judgment but a prayer made in view of it, addressed to the Judge himself, and the name by which it addresses him is \latin{Jesu pie}.

Alphonsus Liguori, writing his considerations on the particular judgment, records that ``hearing the \emph{Dies Ir\ae} sung, and reflecting on the terror of the soul when she shall be presented before the tribunal of Jesus Christ, the venerable P.~Juvenal Ancina took, and afterwards executed, the resolution of forsaking the world.''\endnote{Alphonsus Liguori, \emph{Preparation for Death; or, Considerations on the Eternal Maxims}, Consideration~XXIV, first point (Dublin: James Duffy, 1843), pp.~289--290.} That is what the Church intends the sequence to do. It is not sung over the dead in order to frighten the living for their own sake. It is sung so that the living, praying for the dead, may learn to pray for themselves while there is time.

\subsection{Twice before the Judge}

The tradition the sequence sings presents the Christian standing before Christ the Judge twice. The first time is at the moment of death, alone, when the soul leaves the body and its whole life is set at once in the light of Christ: this the Church calls the particular judgment. The second is at his coming in glory, when all the dead shall rise in their own bodies and the whole human race shall stand together before the seat of his majesty: this is the general or last judgment. Between them lies, for many of the saved, the purification the Church calls purgatory, and over them both lies the charity by which the living help the dead with prayer, alms, and above all the Sacrifice of the altar.

These are not two stories about the same event, nor a provisional hearing followed by the real one. The sentence passed at death is final; the judgment at the end of the world is no less real. The Church holds both with the same firmness and joins them, in the acts that define them, with a single word: \latin{nihilominus}, ``nonetheless.'' To meditate on judgment in the last weeks of the year is to learn why she holds both, and what each asks of a Christian who must one day stand before the Son of Man and does not know the day or the hour.
